This study formulated tangent-only perturbation assessment as a selected
finite-response problem rather than a matrix-magnitude test. Exact algebra
separates the represented perturbation, worst-direction amplification, and
realized forcing direction, while coordinate covariance prevents a
same-numeric diagonal from being mistaken for the same operator after
rescaling.

The main computational result is a pair of a posteriori finite-precision error
enclosures for two algebraically equivalent but arithmetically distinct
operation graphs. The direct enclosure propagates two parent-solve errors and
subtraction error; the response enclosure propagates parent error,
forcing-construction defect, and response-solve residual. Each route covered
300/300 development and 300/300 newly generated holdout systems in both
reference contracts. An independent Decimal implementation confirmed a
stratified 72-lane subset to $6.96\times10^{-153}$ normalized agreement.

The analysis also explains why coverage can coexist with conservative bounds.
Frobenius majorization contributed median factors near unity, whereas
directional-alignment inflation contributed factors of $10^3$--$10^4$.
The response graph often reduced both arithmetic error and enclosure size, but
the binary64 and holdout results ruled out a universal route ranking.

Application to 196- and 900-dimensional coupled finite-element tangents showed
that the full coordinate--response--enclosure chain can qualify archived
application operators. The resulting implementation is a reproducible
reference benchmark for future sparse estimators, not a replacement for their
independent coverage, tightness, and cost validation.
